\begin{abstract}
Vertical ancient Chinese text recognition is a fundamental step in digitizing historical documents. Compared with regular scene text or modern printed text, ancient-book images often contain paper degradation, ink diffusion, low contrast, broken strokes, historical glyph variants, and dataset-specific character inventories. The characters are arranged vertically, and their spacing and local spatial organization are less regular than those in modern printed lines. Existing segmentation-free recognizers can be trained with line-level transcriptions, but they do not explicitly use character-level spatial organization; detection-style recognizers can model character positions, but relying on detected boxes during inference introduces thresholding, filtering, and ordering errors.

We propose SAQT, a structure-aware query learning method for vertical ancient Chinese text recognition. SAQT uses detection-style queries as an intermediate representation between character-localization information and line-level sequence recognition, rather than as final detection outputs. During character-localization information learning, character boxes supervise the queries to capture character-level visual evidence, character positions, and vertical spatial organization. During recognition, query outputs are sorted by vertical position and converted into a CTC sequence trained with line-level transcriptions. To handle character-set mismatch across ancient-text datasets, SAQT further introduces charset-aware classifier adaptation, which inherits classifier parameters for shared characters and initializes rows for newly introduced characters. Experiments on MTHv2 and HDRC show that SAQT obtains higher AR/CR than the vertical-input-adapted scene text recognition baselines considered in this paper; with the development-set-fixed decoding protocol, it reaches 96.75/97.00 and 93.44/94.36 AR/CR, respectively.
\keywords{Ancient text recognition \and Historical document analysis \and Vertical text recognition \and Query-based recognition \and CTC}
\end{abstract}
